\section{Experiment Background}

This project studies a small product-level retrieval problem built from a whisky catalogue. Each product is represented by structured text extracted from the catalogue and by a cropped bottle image. The goal is not to build a full recommendation system with user profiles or purchase history, but to test whether vector embeddings can support similarity search over whisky products using different forms of user input.

We focus on three primary retrieval settings. The first is text-only retrieval, where a natural language query is matched against product text embeddings. The second is image-only retrieval, where either a clean bottle crop or a whole movie scene image is matched against product image embeddings. The third is native multimodal retrieval, where a whole scene image plus the question ``What is the whisky in this image?'' is matched against product embeddings built from both catalogue text and product images.

Beyond these primary settings, we also include cross-representation ablations. In these ablations, a single-modality query is searched against the multimodal product index: text queries are compared with product text-image embeddings, and image queries are compared with product text-image embeddings. These runs test whether adding product-side visual or textual context helps retrieval even when the query itself contains only one modality.

The main question is whether multimodal embeddings provide useful retrieval behaviour beyond modality-matched baselines. We also want to understand the effect of query cleanliness: a cropped bottle image should be easier to match than a full scene image, while a full scene plus text question tests whether the embedding model can use both visual and textual intent. The experiment therefore compares text, image, and image-text retrieval under the same catalogue, product labels, and ranking procedure, and uses ablations to separate the effect of query modality from the effect of product representation.
